\documentclass[10pt,a4paper]{amsart}
\usepackage[margin=19mm]{geometry}
\usepackage{amsmath,amssymb,mathtools}
\usepackage[scr=rsfs,cal=euler]{mathalfa}
\usepackage{xfrac}
\usepackage[unicode,hidelinks,bookmarks=false]{hyperref}
\newcommand{\ie}{i.e.\ }
\newcommand{\cf}{cf.\ }
\newcommand{\resp}{respectively\ }
\newcommand{\Subsection}{\subsection}
\providecommand{\nobreakdash}{\nobreak\hbox{-}}
\setlength{\emergencystretch}{2em}
\multlinegap0pt
\usepackage{bm}
\usepackage{bbm}
\usepackage{dsfont}
\usepackage{xspace}
\usepackage{cancel}
\usepackage{mathtools}
\usepackage{cleveref}
\usepackage{amsthm}
\makeatletter
\@ifundefined{theorem}{\newtheorem{theorem}{Theorem}}{}
\@ifundefined{assumption}{\newtheorem{assumption}[theorem]{Assumption}}{}
\@ifundefined{lemma}{\newtheorem{lemma}[theorem]{Lemma}}{}
\makeatother
\newcommand{\X}{{\mathcal{X}}}
\title[A Modern Theory of Cross-Validation through the Lens of Stab]{A Modern Theory of Cross-Validation through the Lens of Stability}
\author{Jing Lei}
\date{Source: arXiv:2505.23592v3 (CC BY 4.0)}
\begin{document}
\maketitle

\noindent\emph{Source and licence.} A Modern Theory of Cross-Validation through the Lens of Stability. Version arXiv:2505.23592v3 (CC BY 4.0), distributed under the Creative Commons Attribution 4.0 International licence, as recorded in the arXiv licence metadata for this exact version and shown on the abstract page https://arxiv.org/abs/2505.23592v3. Full rights notice: licenses/arxiv-2505.23592-CC-BY-4.0.txt.\par\smallskip

\noindent\textit{Editorial scope.} Counted unit: the Theorem whose source label is \texttt{thm:6.2-second-order-stab-sgd} in the article (source0.tex lines 3203--3207, source label \texttt{thm:6.2-second-order-stab-sgd}), delivered together with the article's own complete original proof of it (source0.tex lines 3388--3532, one \texttt{proof} environment of 145 lines). No mathematics is written, repaired or re-derived: statement and proof are the article's own text line for line. The proof is detached from the statement in the source: the article states the result at lines 3203--3207 and prints its proof at lines 3388--3532, and the proof environment's own header names the statement, so the association is the article's own. Required context retained uncounted: asm:6.1-smooth\_convex (lines 3061--3068); thm:6.1-first-order-stab-sgd (lines 3096--3100); lem:6.1-first-order-error (lines 3110--3115); asm:6.2-smooth-hessian (lines 3189--3195). Every definition, display and label of those lines is retained unchanged. Omitted from this package and not redistributed: 1--3060, 3069--3095, 3101--3109, 3116--3188, 3196--3202, 3208--3387, 3533--3778. Nothing inside a retained excerpt is omitted or altered: every retained body is delivered exactly as the article's lines are written, so no omission receipt of any kind accompanies this package. Citations: the retained excerpts carry no citation command at all, so no bibliography entry of the article is transcribed and no reference is redistributed. Reference adaptation: none was needed and none was made; every reference inside the delivered unit resolves inside it, so no unresolved reference or citation is produced. Printed slips kept literal: no printed word, symbol, hypothesis or number of the counted statement or of its proof was altered. Layout and class: the article's own class is \texttt{article} with packages including geometry, parskip, amsmath, amsthm, amssymb, bbm, bm, mathtools, cases, graphicx, tabularx, microtype, enumitem, dsfont; the wrapper is the lane's pinned \texttt{amsart} editorial wrapper at 10pt on A4, which restates the theorem-like environments the retained text needs and adds the article's own macro definitions that the retained text needs (\texttt{\textbackslash X}), with extra wrapper packages (bm, bbm, dsfont, xspace, cancel, mathtools, cleveref). The delivered numbering is the unit's own. Assets: no retained span contains a figure environment, a graphics-inclusion command, a PDF-page inclusion, a table or a TikZ picture; no image file is copied into this package, the package declares no asset dependency and no image file is redistributed with it. Licence: version arXiv:2505.23592v3 of this article is distributed under the Creative Commons Attribution 4.0 International licence, as recorded in the arXiv licence metadata for this exact version and shown on the abstract page https://arxiv.org/abs/2505.23592v3. A full rights notice is delivered at licenses/arxiv-2505.23592-CC-BY-4.0.txt. Characters: every retained excerpt is pure ASCII, so the delivered engine, XeLaTeX with its default Latin Modern text font, reports no missing character.
\begin{assumption}
    [Smoothness and convexity of objective function]\label{asm:6.1-smooth_convex}
    There exists positive constants $(\gamma, C_0, \beta)$ such that for all $x\in\X$
    \begin{enumerate}
        \item [(a)] $\varphi(\cdot;x)$ is $\gamma$-strongly convex.
        \item [(b)] $\varphi(\cdot;x)$ is $C_0$-Lipschitz
        \item [(c)] $\dot\varphi(\cdot;x)$ is $\beta$-Lipschitz        \end{enumerate}
\end{assumption}

\begin{theorem}[First-order stability of SGD]\label{thm:6.1-first-order-stab-sgd}
    Under \Cref{asm:6.1-smooth_convex}, if the learning rate satisfies $\alpha_t=\beta^{-1}t^{-a}$ for an $a\in(0,1)$ and
    $$\frac{\gamma}{\beta+\gamma}\ge \frac{a(1-a)}{1-2^{-(1-a)}}\frac{\log n}{n^{1-a}}\,,$$
    then for any $i\in [n]$ we have $$\left\|\nabla_i \hat f_n\right\|\le \frac{2^{1+a}C_0}{\beta}n^{-a}\,.$$
\end{theorem}

\begin{lemma}\label{lem:6.1-first-order-error}
Under \Cref{asm:6.1-smooth_convex} and the conditions in \Cref{thm:6.1-first-order-stab-sgd}, we have for all $t\ge i$
$$
\left\|\nabla_i\hat f_t\right\|\le\frac{2C_0}{\beta}i^{-a}\exp\left[-\frac{\gamma}{(1-a)(\beta+\gamma)}\left\{(t+1)^{1-a}-(i+1)^{1-a}\right\}\right]\,.
$$
\end{lemma}

\begin{assumption}\label{asm:6.2-smooth-hessian}
 The Hessian matrix of $\varphi(\cdot;x)$, denoted $\ddot\varphi(\cdot;x)$, is $C_2$-Lipschitz for a constant $C_2$ and all $x$:
 $$
 \|\ddot\varphi(f;x)-\ddot\varphi(f';x)\|_{\rm op}\le C_2\|f-f'\|
 $$
 for all $f\,,f'\in \mathcal F$ and $\|\cdot\|_{\rm op}$ denotes the operator norm.
\end{assumption}

\begin{theorem}[Second-Order Stability of SGD]\label{thm:6.2-second-order-stab-sgd}
    Under the same assumptions as in \Cref{thm:6.1-first-order-stab-sgd}, assume in addition that $\varphi$ satisfies \Cref{asm:6.2-smooth-hessian} such that $\alpha_t\gamma\le 1$ for all $t$, then
    $$\|\nabla_i\nabla_j\hat f\|\le cn^{-2a}\log n\,,$$
    for a constant $c$ depending on $(a,\beta,\gamma,C_0,C_2)$.
\end{theorem}

\begin{proof}[Proof of \Cref{thm:6.2-second-order-stab-sgd}]
     Assume $i<j$. Let $\hat f_t^{i}$ be the estimate at time $t$ with $X_i'$ as the $i$th sample point. Define $\hat f_t^{ij}$ correspondingly.

     Then
\begin{align}
\left\|\nabla_i\nabla_j\hat f_{j}\right\|= & \left\|\hat f_{j} - \hat f^i_{j} - \hat f^j_{j}+\hat f^{ij}_{j}\right\|\nonumber\\
=&\big\|\hat f_{j-1}-\alpha_{j}\dot\varphi(\hat f_{j-1};X_j) -
\hat f^i_{j-1}+\alpha_{j}\dot\varphi(\hat f^i_{j-1};X_j) \nonumber\\
&-\hat f_{j-1}+\alpha_{j}\dot\varphi(\hat f_{j-1};X_j') +
\hat f^i_{j-1}-\alpha_{j}\dot\varphi(\hat f^i_{j-1};X_j')\big\| \nonumber\\
=&\alpha_{j}\left\|\dot\varphi(\hat f_{j-1};X_j)-\dot\varphi(\hat f^i_{j-1};X_j)-\dot\varphi(\hat f_{j-1};X_j')+\dot\varphi(\hat f^i_{j-1};X_j')\right\|\nonumber\\
\le & 2\beta\alpha_{j}  \|\nabla_i\hat f_{j-1}\|\,,\label{eq:6.2-eta_j}\,,
\end{align}
where the last inequality follows from the Lipschitz property of $\dot\varphi(\cdot;x)$.

For $t>j$,
\begin{align}
  \nabla_i\nabla_j\hat f_{t}= & \hat f_{t} - \hat f^i_{t} - \hat f^j_{t}+\hat f^{ij}_{t}\nonumber\\
=&\hat f_{t-1}-\alpha_{t}\dot\varphi(\hat f_{t-1};X_t) -
\hat f^i_{t-1}+\alpha_{t}\dot\varphi(\hat f^i_{t-1};X_t) \nonumber\\
&-\hat f^j_{t-1}+\alpha_{t}\dot\varphi(\hat f^j_{t-1};X_t) +
\hat f^{ij}_{t-1}-\alpha_{t}\dot\varphi(\hat f^{ij}_{t-1};X_t) \nonumber\\
=&\nabla_i\nabla_j \hat f_{t-1} - \alpha_t\left[\dot\varphi(\hat f_{t-1};X_t)-\dot\varphi(\hat f_{t-1}^i;X_t)-\dot\varphi(\hat f_{t-1}^j;X_t)+\dot\varphi(\hat f_{t-1}^{ij};X_t)\right]\,.\label{eq:6.2-nabla2-f-eq1}
\end{align}
Using the mean value theorem, there exist $\tilde f_{t-1}\in[\hat f_{t-1},\hat f_{t-1}^i]$ and $\tilde f_{t-1}'\in[\hat f_{t-1}^j,\hat f_{t-1}^{ij}]$ such that
\begin{align*}
&  \dot\varphi(\hat f_{t-1};X_t)-\dot\varphi(\hat f_{t-1}^i;X_t)-\dot\varphi(\hat f_{t-1}^j;X_t)+\dot\varphi(\hat f_{t-1}^{ij};X_t)\\
=& \ddot\varphi(\tilde  f_{t-1};X_t)(\hat f_{t-1}-\hat f_{t-1}^i)-\ddot\varphi(\tilde f_{t-1}';X_t)(\hat f_{t-1}^j-\hat f_{t-1}^{ij})\\
=&\ddot\varphi(\tilde  f_{t-1};X_t) \nabla_i\nabla_j\hat f_{t-1}+\left[\ddot\varphi(\tilde f_{t-1};X_t)-\ddot\varphi(\tilde f_{t-1}';X_t)\right](\hat f_{t-1}^j-\hat f_{t-1}^{ij})\,.
\end{align*}

Plugging this into the RHS of \eqref{eq:6.2-nabla2-f-eq1}, we obtain
\begin{align*}
  \nabla_i\nabla_j\hat f_t = \left[I-\alpha_t \ddot\varphi(\tilde  f_{t-1};X_t)\right] \nabla_i\nabla_j\hat f_{t-1}+\alpha_t\left[\ddot\varphi(\tilde f_{t-1};X_t)-\ddot\varphi(\tilde f_{t-1}';X_t)\right](\hat f_{t-1}^j-\hat f_{t-1}^{ij})\,.
\end{align*}


Let $\eta_t=\|\nabla_i\nabla_j\hat f_t\|$, we get
\begin{equation}
\eta_t\le (1-\alpha_t\gamma)\eta_{t-1}+2C_2 \alpha_t\|\tilde  f_{t-1}-\tilde f_{t-1}'\|\cdot\|\hat f_{t-1}^j-\hat f_{t-1}^{ij}\|\,,\label{eq:6.2-eta-rec-1}
\end{equation}
where the second term on the RHS uses Lipschitz property of $\ddot\varphi$ (with Lipschitz constant $C_2$).

To control $\tilde f_{t-1}-\tilde f_{t-1}'$, apply \Cref{lem:6.1-first-order-error} directly to $\hat f_{t-1}-\hat f_{t-1}^j$ and $\hat f_{t-1}^i-\hat f_{t-1}^{ij}$, and then use triangle inequality on $\hat f_{t-1}-\hat f_{t-1}^{ij}=(\hat f_{t-1}-\hat f_{t-1}^j)+(\hat f_{t-1}^j-\hat f_{t-1}^{ij})$ and $\hat f^i_{t-1}-\hat f_{t-1}^{j} = (\hat f_{t-1}^i-\hat f_{t-1})+(\hat f_{t-1}-\hat f_{t-1}^j)$, then we have,
\begin{align*}
 & \|\tilde f_{t-1}-\tilde f_{t-1}'\|\\
 \le & \max\left\{\|\hat f_{t-1}-\hat f_{t-1}^j\|\,,~\|\hat f_{t-1}-\hat f_{t-1}^{ij}\|\,,~\|\hat f_{t-1}^i-\hat f_{t-1}^j\|\,,~\|\hat f_{t-1}^i-\hat f_{t-1}^{ij}\|\right\}\\
\le 
% & \frac{2C_0}{\beta}j^{-a}\exp\left[-c_{a,\beta,\gamma}((t+1)^{1-a}-(j+1)^{1-a})\right]+
% \frac{2C_0}{\beta}i^{-a}\exp\left[-c_{a,\beta,\gamma}((t+1)^{1-a}-(i+1)^{1-a})\right]\\
% \coloneqq 
& \frac{2C_0}{\beta}\left(\delta_{t-1,i}(c)i^{-a}+\delta_{t-1,j}(c)j^{-a}\right)\,,
\end{align*}
where $c=\frac{\gamma}{(1-a)(\beta+\gamma)}$ and
$$
\delta_{t,j}(c) = \exp\left[-c((t+1)^{1-a}-(j+1)^{1-a})\right]\,.
$$


For $t\ge j+1$, define 
\begin{align*}
\varrho_{i,j,t}=&\frac{8C_2C_0^2}{\beta^3}t^{-a}\left[\delta_{t-1,i}(c)i^{-a}+\delta_{t-1,j}(c)j^{-a}\right]\delta_{t-1,i}(c)i^{-a}\,,
\end{align*}
then \eqref{eq:6.2-eta-rec-1} implies, for $t\ge j+1$,
\begin{equation}\label{eq:6.2-eta-rec-2}
\eta_{t}\le (1-\alpha_{t}\gamma)\eta_{t-1}+\varrho_{i,j,t}\,.
\end{equation}


We now seek to upper-bound $\eta_n$ using the recursion \eqref{eq:6.2-eta-rec-1}. To do so, let $c'=\gamma/[(1-a)\beta]\ge c$. Thus
\begin{align}
\eta_n \le & (1-\alpha_{n}\gamma)\eta_{n-1}+\varrho_{i,j,n}\nonumber\\
\le & (1-\alpha_{n}\gamma)\left[(1-\alpha_{n-1}\gamma)\eta_{n-2}+\varrho_{i,j,n-1}\right]+\varrho_{i,j,n}\nonumber\\
% \le & \cdots\\
\le & \left[\prod_{k=j+1}^n(1-\alpha_k\gamma)\right] \eta_j + \sum_{k=j+1}^{n} \left[\prod_{l=k+1}^n(1-\alpha_l\gamma)\right] \varrho_{i,j,k}\nonumber\\
\le & \delta_{n,j}(c')\eta_j+
\sum_{k=j+1}^{n} \delta_{n,k}(c')\varrho_{i,j,k}\nonumber\\
=& \underbrace{\delta_{n,j}(c')\eta_j}_{(III)}+
\underbrace{\frac{8C_2C_0^2}{\beta^3}\sum_{k=j+1}^{n} \delta_{n,k}(c')\delta_{k-1,i}^2(c) k^{-a}i^{-2a}}_{(II)}\nonumber\\
&~~~+\underbrace{\frac{8C_2C_0^2}{\beta^3}\sum_{k=j+1}^{n} \delta_{n,k}(c')\delta_{k-1,i}(c)\delta_{k-1,j}(c) k^{-a}i^{-a}j^{-a}}_{(I)}\,,\label{eq:6.2-eta-rec-3}
\end{align}
where the last inequality appeals to the fact that for positive integers $j<t$, $\prod_{k=j+1}^t (1-\alpha_k \gamma)\le \exp\left[-\frac{\gamma}{(1-a)\beta}\left((t+1)^{1-a}-(j+1)^{1-a}\right)\right]$.

\textbf{Term (I):} Using the fact that $\delta_{k-1,i}(c)\ge \delta_{k-1,i}(c')$, $\delta_{n,k}(c)\delta_{k-1,i}(c)\le \delta_{n,i}(c)e^{c(2^{1-a}-1)}$, and $\delta_{k-1,j}(c)\le 1$, we have
\begin{align}
    (I)\le & \frac{8C_2C_0^2}{\beta^3}e^{c(2^{1-a}-1)}i^{-a}j^{-a}\delta_{n,i}(c)\sum_{k=j+1}^n k^{-a}\nonumber\\
    \le & \frac{8C_2C_0^2}{(1-a)\beta^3}e^{c(2^{1-a}-1)}i^{-a}j^{-a}\delta_{n,i}(c)(n^{1-a}-j^{1-a})\,.\label{eq:6.2-term-I-init}
\end{align}


Let $c_2=(3-a)/[c(1-a)]$.
When $i\le n-c_2 (n+1)^a \log n$, we have
\begin{align*}
  (n+1)^{1-a}-(i+1)^{1-a} = & (n+1)^{1-a}\left[1-\left(1-\frac{n-i}{n+1}\right)^{1-a}\right]\\
  \ge & (n+1)^{1-a}(1-a)\frac{n-i}{n+1}\\
  = & c_2(1-a) \log n\,,
\end{align*}
where the second lines uses $(1-x)^b\le 1-bx$ for $x\in[0,1)$ and $b\in(0,1)$. Thus 
\eqref{eq:6.2-term-I-init} implies
\begin{align}
   (I)\le & \frac{8 C_2C_0^2 }{(1-a)\beta^3}e^{c(2^{1-a}-1)}i^{-a}j^{-a}n^{-(3-a)}(n^{1-a}-j^{1-a})\nonumber \\
    \le & \frac{8 C_2C_0^2 }{(1-a)\beta^3}e^{c(2^{1-a}-1)} n^{-2}\,.  \label{eq:6.2-term-I-small-i}
\end{align}

When $i\ge n-c_2 (n+1)^a\log n$.  Suppose $n$ is large enough that $n-c_2(n+1)^a\log n\ge n/2$. Use the following upper bounds in the RHS of \eqref{eq:6.2-term-I-init}: 
$k^{-a}\le j^{-a}\le i^{-a}\le 2^{a} n^{-a}$, $\delta_{n,i}(c)\le 1$, we have
\begin{align}
   (I)\le  & \frac{8C_2 C_0^2 }{(1-a)\beta^3}e^{c(2^{1-a}-1)}2^{3a}n^{-3a}(n-j)\nonumber\\
    \le & \frac{8C_2 C_0^2 }{(1-a)\beta^3}e^{c(2^{1-a}-1)}2^{3a}n^{-3a}(n-i)\nonumber\\
    \le & \frac{8C_2 C_0^2 }{(1-a)\beta^3}e^{c(2^{1-a}-1)}2^{3a}n^{-3a}c_2(n+1)^a\log n\nonumber\\
    \le & \frac{8C_2 C_0^2 }{(1-a)\beta^3}e^{c(2^{1-a}-1)}2^{3a+1}c_2 n^{-2a}\log n\,.\label{eq:6.2-term-I-i-large}
\end{align}
Combining \eqref{eq:6.2-term-I-small-i} and \eqref{eq:6.2-term-I-i-large}, we conclude that
\begin{equation}
    \label{eq:6.2-term-I-final}
    (I)\le c(a,\beta,\gamma,C_0,C_2) n^{-2a}\log n\,,
\end{equation}
for $n$ large enough.

\textbf{Term (II):}
The second term in  \eqref{eq:6.2-eta-rec-3} is upper bounded by, for large enough $n$,
\begin{align}
  (II)\le &\frac{8 C_2 C_0^2 }{\beta^3}i^{-2a} \sum_{k=j+1}^n k^{-a} \delta_{n,k}(c)
 \delta_{k-1,i}(c)\nonumber\\
 \le & \frac{8 C_0^2 C_2}{(1-a)\beta^3} e^{c(2^{1-a}-1)}i^{-2a} \delta_{n,i}(c)(n^{1-a}-j^{1-a}) \nonumber\\
 \le & c(a,\beta,\gamma,C_0,C_2) n^{-2a}\log n\,,\label{eq:6.2-term-II-final}
\end{align}
where the first inequality follows from $\delta_{k-1,i}(c)\le 1$, the second and third follow from the same argument as in term (I).

\textbf{Term (III):}
\begin{align}
(III)=  &\delta_{n,j}(c')\eta_j\nonumber\\
\le & 2\beta\alpha_j \|\nabla_i\hat f_{j-1}\|\delta_{n,j}(c')\nonumber\\
\le & \frac{4C_0}{\beta}j^{-a}i^{-a}\delta_{j-1,i}(c')\delta_{n,j}(c')\nonumber\\
\le & \frac{4C_0}{\beta}e^{c(2^{1-a}-1)}j^{-a}i^{-a}\delta_{n,i}(c)\nonumber\\
\le & \frac{4C_0}{\beta}e^{c(2^{1-a}-1)}i^{-2a}\delta_{n,i}(c)\,,\label{eq:6.2-term-III-init}
\end{align}
where the first inequality follows from \eqref{eq:6.2-eta_j}, the second from \Cref{lem:6.1-first-order-error}, the third from $\delta_{j-1,i}(c)\ge \delta_{j-1,i}(c')$ and the same argument used in the treatment of term (I), and the fourth one from the fact that $j>i$. 

Finally, a similar argument partitioning $i$ according to $i\le n/2$ and $i\ge n/2$ as in the proof of \Cref{thm:6.1-first-order-stab-sgd} leads to the conclusion
\begin{equation}\label{eq:6.2-term-III-final}
	(III)\le c(a,\beta,\gamma,C_0,C_2) n^{-2a}\,.
\end{equation}
The proof concludes by plugging \eqref{eq:6.2-term-I-final}, \eqref{eq:6.2-term-II-final}, and \eqref{eq:6.2-term-III-final} into \eqref{eq:6.2-eta-rec-3}.
\end{proof}

\end{document}
